% ============================================================================
% HISTORICAL / DO NOT PUBLISH
% This superseded Q=3360 draft is retained only as a development record.
% The current manuscript is paper/erdos302_two_sided.tex.
% ============================================================================
\documentclass[11pt]{article}
\usepackage[a4paper,margin=30mm]{geometry}
\usepackage{amsmath,amssymb,amsthm}
\usepackage[hidelinks]{hyperref}

\newtheorem{theorem}{Theorem}
\newtheorem{lemma}{Lemma}
\newcommand{\N}{\mathbb N}

\title{A computer-assisted upper bound for Erd\H{o}s problem 302}
\author{Author name to be supplied before publication}
\date{August 2026}

\begin{document}
\thispagestyle{empty}
\vspace*{\fill}
\begin{center}
\fbox{\begin{minipage}{0.92\linewidth}
\centering\bfseries
HISTORICAL / DO NOT PUBLISH

\medskip
This superseded \(Q=3360\) draft is retained only as a development record.
The current manuscript is \texttt{paper/erdos302\_two\_sided.tex} and proves
the stronger certified upper bound
\(140803024/163562355\), together with the lower-bound progress.
\end{minipage}}
\end{center}
\vspace*{\fill}
\clearpage
\setcounter{page}{1}

\maketitle

\begin{abstract}
Let $f(N)$ be the largest size of a set $A\subseteq\{1,\ldots,N\}$
containing no three distinct elements $a,b,c$ with
$1/a=1/b+1/c$.  We give a computer-assisted proof that
\[
  \limsup_{N\to\infty}\frac{f(N)}N\leq \frac{5273}{6048}
  =0.8718584656\ldots .
\]
The finite lemma is proved by a dependency-free exhaustive verifier and is
independently cross-checked by a mixed-integer linear program.  Only the
elementary arithmetic layer is currently checked in Lean 4; this note does
not claim an end-to-end Lean formalization.
\end{abstract}

\section{The finite hypergraphs}
Put $q=3360=2^5\cdot3\cdot5\cdot7$ and
\[
 D=\operatorname{Div}(q)\setminus\{1\}.
\]
Thus $|D|=47$.  For $t>0$, let $H_t$ be the three-uniform hypergraph with
vertex set $D_{\leq t}=\{d\in D:d\leq t\}$ and edges the unordered triples
$\{a,b,c\}$ of distinct vertices satisfying $bc=a(b+c)$ (with $a<b<c$ used
only to serialize an edge once).

\begin{lemma}[finite certificate]\label{lem:certificate}
For
\begin{align*}
(t_1,\ldots,t_{21})={}&(6,12,24,30,40,42,56,84,96,105,120,140,\\
&210,224,240,336,420,560,1120,1680,3360),
\end{align*}
the vertex-cover number of $H_{t_r}$ is exactly $r$.
\end{lemma}

The theorem below needs only the lower bounds $\tau(H_{t_r})\geq r$.
The accompanying exact verifier regenerates $D$ and all 146 edges directly
from the displayed definitions.  To disprove a cover of size at most $r-1$,
it chooses an uncovered edge and branches over the three necessary choices
of a vertex from that edge.  A branch succeeds only after all edges have been
covered.  A branch is rejected when its budget is exhausted, or when either
of two rigorous lower bounds exceeds the remaining budget: a packing of
pairwise disjoint edges, and $\lceil |E|/\Delta\rceil$, where $\Delta$ is the
maximum remaining degree.  These rejections cannot remove a feasible cover.
Induction on the remaining budget therefore proves that the search explores
all possible covers.  Memoization merely reuses the result for an identical
pair (remaining edges, budget).

For the reverse inequalities, the verifier contains an explicit set of $r$
vertices for every $t_r$ and checks by integer bit operations that it meets
every edge.  Consequently the equality in Lemma~\ref{lem:certificate} does
not depend on the floating-point MILP.  The recorded SHA-256 digest fixes the
canonical edge serialization and detects accidental data changes; it is not
itself a proof of any mathematical claim.

For positive integers, the edge equation has precisely the intended meaning:
\[
 bc=a(b+c)
 \quad\Longleftrightarrow\quad
 \frac1a=\frac1b+\frac1c,
\]
because multiplying the latter equality by the nonzero integer $abc$ gives
the former (and division by $abc$ proves the converse).  Scaling all three
denominators by the same positive integer preserves the equality.

\section{Disjoint multiplier blocks}
For $L\geq1$ define the multiplier set
\[
 \mathcal M_L=\left\{
 2^{6\alpha}3^{2\beta}5^{2\gamma}7^{2\delta}u:
 0\leq\alpha,\beta,\gamma,\delta<L,\quad (u,210)=1
 \right\}.
\]
Although $\mathcal M_L$ is infinite, every truncated set
$\mathcal M_L(X)=\mathcal M_L\cap[1,X]$ is finite; all sums below use these
finite truncations.

\begin{lemma}\label{lem:disjoint}
The blocks $mD$, for $m\in\mathcal M_L$, are pairwise disjoint.
\end{lemma}
\begin{proof}
Suppose $md=m'd'$ with $m,m'\in\mathcal M_L$ and $d,d'\mid q$.
Modulo $6$, the $2$-adic valuation of the equality gives
$v_2(d)=v_2(d')$, since both lie in $[0,5]$.  Modulo $2$, the same argument
at $3,5,7$ gives equality of the corresponding valuations, since each lies
in $[0,1]$.  A divisor of $q$ is determined by these four valuations, so
$d=d'$, and cancellation gives $m=m'$.
\end{proof}

\begin{lemma}\label{lem:density}
For fixed $L$,
\[
 |\mathcal M_L(X)|=\rho_LX+O_L(1),
 \quad
 \rho_L=\prod_{(p,k)\in\{(2,6),(3,2),(5,2),(7,2)\}}
 \frac{1-p^{-1}}{1-p^{-k}}(1-p^{-kL}),
\]
and $\lim_{L\to\infty}\rho_L=5/18$.
\end{lemma}
\begin{proof}
For a fixed quadruple $(\alpha,\beta,\gamma,\delta)$, inclusion--exclusion
over the four primes shows that the count of eligible $u$ up to $Y$ is
$Y\prod_{p\mid210}(1-1/p)+O(1)$.  There are only $L^4$ quadruples, so summing
these estimates is legitimate and yields the displayed product of four
finite geometric series.  Letting $L$ tend to infinity gives
\[
 \prod_{(p,k)}\frac{1-p^{-1}}{1-p^{-k}}
 =\frac{1/2}{1-2^{-6}}\frac{2/3}{1-3^{-2}}
  \frac{4/5}{1-5^{-2}}\frac{6/7}{1-7^{-2}}=\frac5{18}.
\]
\end{proof}

\section{Counting forced omissions}
Let $A\subseteq[1,N]$ contain no forbidden reciprocal triple.  For each
$m\in\mathcal M_L$, the omitted vertices in the finite block
$mD\cap[1,N]$ form a vertex cover of the relevant prefix hypergraph: an
uncovered edge would scale to a forbidden triple in $A$.  If
$m\leq N/t_r$, Lemma~\ref{lem:certificate} forces at least $r$ omissions in
$mD_{\leq t_r}$.  Hence the number of omissions in the whole block is at
least the number of indices $r$ for which $m\leq N/t_r$.  Summing this
statement and using Lemma~\ref{lem:disjoint} gives the finite inequality
\[
 N-|A|\geq\sum_{r=1}^{21}
 \left|\mathcal M_L\!\left(\left\lfloor\frac N{t_r}\right\rfloor\right)\right|.
\]
By Lemma~\ref{lem:density}, for fixed $L$, division by $N$ followed by
$N\to\infty$ yields
\[
 \liminf_{N\to\infty}\frac{N-f(N)}N
 \geq \rho_L\sum_{r=1}^{21}\frac1{t_r}.
\]
Only after this first limit do we let $L\to\infty$.  Direct exact rational
arithmetic gives
\[
 \sum_{r=1}^{21}\frac1{t_r}=\frac{155}{336},\qquad
 \frac5{18}\frac{155}{336}=\frac{775}{6048}.
\]
We have proved the following.

\begin{theorem}
\[
 \limsup_{N\to\infty}\frac{f(N)}N
 \leq1-\frac{775}{6048}=\frac{5273}{6048}.
\]
Equivalently, $f(N)\leq(5273/6048+o(1))N$.
\end{theorem}

\section*{Computational and AI disclosure}
The finite configuration and initial proof strategy were developed with
substantial assistance from OpenAI models.  Codex was used to prepare and
review verification code, Lean infrastructure, and this manuscript.  The
human author must review the argument, replace the author placeholder, select
the final claims, and assume responsibility before publication.  The present
repository establishes a computer-assisted partial result; it does not claim
that Erd\H{o}s problem 302 is solved or fully formalized in Lean.

\end{document}
